% Fragmento público generado desde propietarios íntegros.
% Residencia: 52.7.2.
% Fuente: ../incoming/radion_monografia_20260827/manuscrito/sections/08_complejo_electromagnetismo.tex:59-158
\subsubsection{Operador constitutivo sobre los canales \(90/120\)}

\paragraph{Descomposición en componentes común y orientada}

Definimos
\begin{align}
\mathcal E
&=\log\bigl[(1-q_+^{90})(1-q_-^{90})\bigr]
-\log\bigl[(1-q_+^{120})(1-q_-^{120})\bigr],
\label{eq:Ecommon}\\
\mathcal B
&=\log\frac{1-q_-^{90}}{1-q_+^{90}}
-\log\frac{1-q_-^{120}}{1-q_+^{120}}.
\label{eq:Borient}
\end{align}
Bajo \(C^*\mapsto-C^*\), se intercambian \(q_+\leftrightarrow q_-\). Por
tanto, \(\mathcal E\) es par y \(\mathcal B\) es impar.

Las constitutivas normalizadas son
\begin{equation}
\widehat\mu=e^{\mathcal E+\mathcal B},
\qquad
\widehat\varepsilon=e^{\mathcal E-\mathcal B},
\qquad
\widehat Z=e^{\mathcal B},
\qquad
\widehat c=e^{-\mathcal E}.
\label{eq:constitutives}
\end{equation}
Se verifican exactamente
\begin{equation}
\widehat Z=\sqrt{\frac{\widehat\mu}{\widehat\varepsilon}},
\qquad
\widehat c=(\widehat\mu\widehat\varepsilon)^{-1/2}.
\label{eq:constitutive-identities}
\end{equation}
La conjugación produce
\begin{equation}
C^*\mapsto-C^*:\qquad
\widehat\mu\leftrightarrow\widehat\varepsilon,
\quad
\widehat Z\leftrightarrow\widehat Z^{-1},
\quad
\widehat c\mapsto\widehat c.
\label{eq:EM-conjugation}
\end{equation}

\begin{theorem}[Polaridad constitutiva]
El lector \(90/120\) convierte el modo común y el contraste orientado del
par \((A,C^*)\) en dos constitutivas conjugadas. La inversión de hoja
intercambia permeabilidad y permitividad, invierte impedancia y conserva la
velocidad normalizada.
\end{theorem}

Ésta es la formulación precisa de la intuición \enquote{parte eléctrica y
parte magnética}. No se escribe \(A=E\) ni \(C^*=B\). Se escribe el mapa
\[
(A,C^*)\to(x,y)\to(q_+,q_-)
\to(\mathcal E,\mathcal B)
\to(\widehat\varepsilon,\widehat\mu,\widehat Z,\widehat c).
\]

\paragraph{Realización elástica--rotacional de la polaridad constitutiva}

La misma información puede publicarse en el régimen complejo mediante
\begin{equation}
Z_{\mathrm{em}}=e^{\mathcal E/2}
\left(\cos\frac{\mathcal B}{2}\,I
-\sin\frac{\mathcal B}{2}\,J\right).
\label{eq:Zem}
\end{equation}
El factor \(e^{\mathcal E/2}\) es una dilatación elástica; el factor generado
por \(J\) es una rotación. La escritura compleja no añade una dimensión
imaginaria misteriosa a un mundo real previamente dado: registra dos modos
operacionales que el TRIT ya separó.

\paragraph{Transducción de la orientación en carga efectiva}

La hoja \(\sigma\) del radión fija una orientación antes de introducir una
unidad de carga. La Mecánica Dimensional dispone luego el transductor
\[
e_{\mathrm{int}}=\sqrt{\frac{2\alphah\hfull}{Z_{\mathrm{int}}}},
\qquad
Q=n_Qe_{\mathrm{int}}.
\]
La cadena causal es
\[
\text{hoja}\to\text{orientación}\to\text{constitutiva}
\to\text{transductor dimensional}\to\text{carga física}.
\]
Llamar \emph{radión} a la unidad no identifica prematuramente la orientación
con el culombio; hace visible el lugar donde el signo de carga puede residir.

\begin{narrativebox}
La electricidad y el magnetismo no se pegan al círculo desde fuera. Emergen
cuando el mismo estado se lee a través de producto y cociente: el producto
retiene el modo común, el cociente retiene la orientación. El oscilador LC
realiza el intercambio; el lector \(90/120\) publica las constitutivas; la
carta dimensional asigna unidades.
\end{narrativebox}
